\section*{Problems}

\subsection*{Problem statements}

% Remember
\begin{problema}
	\label{prob:8794f31}
	Let $A, B \subset S$ be two events such that $P(A) = 0.60$, $P(B) = 0.50$, and the joint probability is $P(A \cap B) = 0.30$. Compute directly:
	\begin{enumerate}
		\item $P(A|B)$ and $P(B|A)$
		\item $P(A \cap B')$
		\item $P(A|B')$
	\end{enumerate}
\end{problema}

% Understand
\begin{problema}
	\label{prob:cf929e4}
	An urn contains 5 red balls and 3 blue balls. Two balls are drawn successively without replacement. Let $A$ = \texttt{``the first ball is red''} and $B$ = \texttt{``the second ball is red''}.
	\begin{enumerate}
		\item Compute $P(A)$
		\item Compute $P(B|A)$ and $P(B|A')$
		\item Compute $P(B)$ using the total probability rule, and explain why $P(B)=P(A)$ even though the drawing is without replacement.
	\end{enumerate}
\end{problema}

% Apply
\begin{problema}
	\label{prob:52e63a3}
	In an automated plant, 70\% of the processed parts come from production line A and 30\% from line B. Line A has a defect rate of 5\%, while line B has a defect rate of 10\%. If the quality-control team selects one part at random:
	\begin{enumerate}
		\item What is the total probability that the part is defective?
		\item If the selected part is defective, what is the conditional probability that it came from line A?
	\end{enumerate}
\end{problema}

% Analyze
\begin{problema}
	\label{prob:8a7282b}
	\begin{enumerate}
		\item Formally prove the General Chain Rule (or Multiplication Rule) for a finite sequence of $n$ events with $P(A_1 \cap \dots \cap A_{n-1}) > 0$:
		\begin{align}
			P(A_1 \cap A_2 \cap \dots \cap A_n) = P(A_1)P(A_2|A_1)P(A_3|A_1 \cap A_2) \cdots P(A_n|A_1 \cap \dots \cap A_{n-1}).
		\end{align}
		\item Use this rule to compute the exact probability that, when drawing 4 cards successively and without replacement from a standard 52-card deck, the 4 cards belong to different suits.
	\end{enumerate}
\end{problema}

% Evaluate
\begin{problema}
	\label{prob:c840e8f}
	\textbf{The Monty Hall Paradox:} In a game show, a contestant faces three closed doors ($P_1, P_2, P_3$). Behind exactly one door is the main prize, and behind the other two are goats. The contestant chooses one door at random (suppose, without loss of generality, Door 1). The host, who knows with certainty what is behind each door and is required by the rules always to open a goat door different from the contestant's chosen door (if two options are available, he chooses each with probability $0.5$), opens Door 3 and shows a goat.

	Using the formal axioms of conditional probability, evaluate and prove mathematically why the probability of winning if the contestant switches to Door 2 is exactly $\dfrac{2}{3}$, while if the contestant keeps Door 1 the probability is only $\dfrac{1}{3}$, and explain why the intuition that "two doors remain, so it is 50-50" is incorrect.
\end{problema}

% Create
\begin{problema}
	\label{prob:4241cab}
	Design an original medical-diagnosis problem: a rapid test for detecting a rare disease that affects $1\%$ of the population has sensitivity of $95\%$ (it correctly detects sick people) and specificity of $90\%$ (it correctly gives a negative result for healthy people). Use Bayes' theorem to compute the probability that a person with a positive result actually has the disease, and comment on why this result---known as the \emph{base-rate fallacy}---often surprises people without training in probability.
\end{problema}

\subsection*{Detailed solutions}

% Remember
\begin{solproblema}[prob:8794f31]
	Given $P(A)=0.60$, $P(B)=0.50$, and $P(A \cap B)=0.30$:
	\begin{enumerate}
		\item Direct conditional probabilities:
		\begin{align}
			P(A|B) &= \frac{P(A \cap B)}{P(B)} = \frac{0.30}{0.50} = 0.60 \\
			P(B|A) &= \frac{P(A \cap B)}{P(A)} = \frac{0.30}{0.60} = 0.50
		\end{align}
		(Observe that here $P(A|B)=P(A)$, reflecting independence in this particular exercise).

		\item The probability of $A$ and not $B$ is:
		\begin{align}
			P(A \cap B') = P(A) - P(A \cap B) = 0.60 - 0.30 = 0.30.
		\end{align}

		\item Since $P(B') = 1 - P(B) = 1 - 0.50 = 0.50$:
		\begin{align}
			P(A|B') = \frac{P(A \cap B')}{P(B')} = \frac{0.30}{0.50} = 0.60.
		\end{align}
	\end{enumerate}
\end{solproblema}

% Understand
\begin{solproblema}[prob:cf929e4]
	The urn contains a total of $5 + 3 = 8$ balls.
	\begin{enumerate}
		\item The initial probability for the first ball is $P(A) = \frac{5}{8}$.
		\item If the first ball was red, then 4 red and 3 blue balls remain in the urn ($7$ total): $P(B|A) = \frac{4}{7}$. If the first ball was blue ($A'$), then 5 red and 2 blue balls remain ($7$ total): $P(B|A') = \frac{5}{7}$.
		\item By the total probability rule (partition into $A$ and $A'$):
		\begin{align}
			P(B) &= P(B|A)P(A) + P(B|A')P(A') = \left(\frac{4}{7} \cdot \frac{5}{8}\right) + \left(\frac{5}{7} \cdot \frac{3}{8}\right) = \frac{20}{56} + \frac{15}{56} = \frac{35}{56} = \frac{5}{8}.
		\end{align}
		The fact that $P(B)=P(A)=5/8$ is not accidental: before observing any result, there is no information that distinguishes the second position from the first within the drawing sequence---by symmetry, any ball in the urn has the same marginal probability of ending up as "the red ball in position 2" as "the red ball in position 1"---although the specific result of the first draw does affect the \emph{conditional} probability of the second.
	\end{enumerate}
\end{solproblema}

% Apply
\begin{solproblema}[prob:52e63a3]
	Let $D$ be the event that the part is defective, $A$ the event that it comes from line A, and $B$ the event that it comes from line B.
	\begin{enumerate}
		\item By the total probability rule:
		\begin{align}
			P(D) &= P(D|A)P(A) + P(D|B)P(B) = (0.05)(0.70) + (0.10)(0.30) = 0.035 + 0.030 = 0.065 \text{ (6.5\%).}
		\end{align}
		\item By Bayes' theorem, given that $D$ occurred:
		\begin{align}
			P(A|D) = \frac{P(D|A)P(A)}{P(D)} = \frac{0.035}{0.065} = \frac{7}{13} \approx 0.5385 \text{ (53.85\%).}
		\end{align}
	\end{enumerate}
\end{solproblema}

% Analyze
\begin{solproblema}[prob:8a7282b]
	\begin{enumerate}
		\item \textbf{Proof by induction of the Chain Rule:} for $n=2$, the formula reduces to $P(A_1 \cap A_2) = P(A_1)P(A_2|A_1)$, which is exactly the definition of conditional probability. Assume as the induction hypothesis that the identity holds for $k$ simultaneous events. For $k+1$ events, group the first $k$ into a single compound event $E_k = A_1 \cap A_2 \cap \dots \cap A_k$:
		\begin{align}
			P(A_1 \cap \dots \cap A_k \cap A_{k+1}) = P(E_k \cap A_{k+1}) = P(E_k) \cdot P(A_{k+1}|E_k).
		\end{align}
		Substituting the expansion of $P(E_k)$ from the induction hypothesis immediately gives the product on the right-hand side with the $(k+1)$-st appended term, proving the induction for all $n$.

		\item \textbf{Drawing 4 cards of different suits:} let $C_i$ be the event that the $i$-th card is of a suit different from all previous cards.
		\begin{itemize}
			\item For $C_1$: the first card can be anything $\implies P(C_1) = 1$.
			\item For $C_2|C_1$: of the 51 remaining cards, 39 belong to the 3 suits not yet chosen $\implies P(C_2|C_1) = \frac{39}{51}$.
			\item For $C_3|C_1 \cap C_2$: 50 cards remain and 2 suits are still unseen ($26$ useful cards) $\implies P(C_3|\dots) = \frac{26}{50}$.
			\item For $C_4|C_1 \cap C_2 \cap C_3$: 49 cards remain and 1 suit is still unseen (13 useful cards) $\implies P(C_4|\dots) = \frac{13}{49}$.
		\end{itemize}
		By the Chain Rule:
		\begin{align}
			P(C_1 \cap C_2 \cap C_3 \cap C_4) = 1 \times \frac{39}{51} \times \frac{26}{50} \times \frac{13}{49} \approx 0.1055 \text{ (10.55\%).}
		\end{align}
	\end{enumerate}
\end{solproblema}

% Evaluate
\begin{solproblema}[prob:c840e8f]
	Formalize the conditional space of the game. Let $C_i$ be the event that the prize is behind Door $i$, with $P(C_1)=P(C_2)=P(C_3)=\frac{1}{3}$; and let $A_3$ be the event that the host opens Door 3 and shows a goat. Evaluate the conditional probabilities of the host's action for each possible prize location:
	\begin{align}
		P(A_3 | C_1) &= \frac{1}{2} \quad \text{(if the prize is behind 1, the host chooses randomly between 2 and 3)} \\
		P(A_3 | C_2) &= 1 \quad \text{(if the prize is behind 2, the host must open 3)} \\
		P(A_3 | C_3) &= 0 \quad \text{(if the prize is behind 3, the host never opens it)}
	\end{align}
	By the total probability rule:
	\begin{align}
		P(A_3) = \left(\frac{1}{2} \cdot \frac{1}{3}\right) + \left(1 \cdot \frac{1}{3}\right) + \left(0 \cdot \frac{1}{3}\right) = \frac{1}{6} + \frac{1}{3} = \frac{1}{2}.
	\end{align}
	Applying Bayes: $P(C_1|A_3) = \dfrac{(1/2)(1/3)}{1/2} = \dfrac{1}{3}$ and $P(C_2|A_3) = \dfrac{(1)(1/3)}{1/2} = \dfrac{2}{3}$.

	The intuition "two doors remain, so it is 50-50" is incorrect because it ignores that the host's action is \textbf{not random or symmetric with respect to the information it provides}: the host uses certain knowledge of the door contents to \emph{always} avoid opening the prize door, which transfers all the probability originally held by the two unchosen doors to the only unchosen door that remains closed. Door 1 never updates its probability ($1/3$ before and after), but Door 2 absorbs the full probability that it previously shared with Door 3 ($1/3+1/3=2/3$).
\end{solproblema}

% Create
\begin{solproblema}[prob:4241cab]
	Let $D$ be the event of being sick, with $P(D)=0.01$ and $P(D')=0.99$. The sensitivity gives $P(+|D)=0.95$, and the specificity of $90\%$ implies a false-positive rate $P(+|D')=1-0.90=0.10$.

	The total probability of a positive result is:
	\begin{align}
		P(+) = P(+|D)P(D) + P(+|D')P(D') = (0.95)(0.01) + (0.10)(0.99) = 0.0095 + 0.099 = 0.1085.
	\end{align}
	By Bayes' theorem:
	\begin{align}
		P(D|+) = \frac{P(+|D)P(D)}{P(+)} = \frac{0.0095}{0.1085} \approx 0.0876 \text{ (8.76\%).}
	\end{align}
	Although the test seems highly reliable (95\% sensitivity, 90\% specificity), only $8.76\%$ of those who receive a positive result actually have the disease. This is surprising because intuition focuses on the accuracy of the test itself, ignoring that the disease is so rare ($1\%$) that the absolute number of false positives among the vast majority of healthy people ($0.10 \times 99\%$ of the population) greatly exceeds the number of true positives among the small sick minority ($0.95\times1\%$).
\end{solproblema}
